% Generated draft: requires mathematical and source review.
\documentclass[11pt]{article}
\usepackage[margin=25mm]{geometry}
\usepackage{fontspec}
\setmainfont{Latin Modern Roman}
\usepackage{amsmath,amssymb,mathtools,unicode-math}
\setmathfont{Latin Modern Math}
\usepackage{enumitem,needspace,xurl,hyperref}
\setlength{\parindent}{0pt}
\setlength{\parskip}{5pt}
\setlength{\emergencystretch}{3em}
\newcommand{\N}{\mathbb{N}}
\newcommand{\Z}{\mathbb{Z}}
\newcommand{\Q}{\mathbb{Q}}
\newcommand{\R}{\mathbb{R}}
\newcommand{\C}{\mathbb{C}}
\newcommand{\F}{\mathbb{F}}
\newcommand{\A}{\mathbb{A}}
\newcommand{\PP}{\mathbb{P}}
\newcommand{\E}{\mathbb{E}}
\newcommand{\eps}{\varepsilon}
\newcommand{\dd}{\,\mathrm{d}}
\begin{document}
\section*{UnsolvedMath ID: 701130}
\textbf{Source problem code:} catalogue entry
\subsection{Definitions and mathematical statement}
Let \(F\) denote R. Thompson's group, in its standard realization as the group under composition of increasing piecewise linear homeomorphisms of \([0,1]\) with finitely many breakpoints, all breakpoints dyadic rational, and all slopes integral powers of \(2\). A dyadic rational is a number of the form \(m/2^n\), where \(m\in\Z\) and \(n\geq 0\). These homeomorphisms fix both endpoints.

Let \(k\) be an arbitrary field. The group algebra \(k[F]\) consists of finite formal sums \(\sum_{g\in F}c_g g\), with \(c_g\in k\), with coefficientwise addition and multiplication induced by the group operation and distributivity.

The question is whether \(k[F]\) satisfies the following right Ore condition, in exactly the formulation supplied by the source:
\[
\forall a,b\in k[F]\;\exists u,v\in k[F]\quad
au=bv\quad\text{and}\quad(u\neq 0\ \lor\ v\neq 0).
\]
There is no restriction in the supplied statement on the characteristic of \(k\). The displayed condition includes the cases where \(a\) or \(b\) is zero, and requires only that at least one of \(u,v\) be nonzero.

For the accompanying qualification, a group \(G\) is amenable if there exists a finitely additive function \(\mu:\mathcal P(G)\to[0,1]\) such that \(\mu(G)=1\) and \(\mu(gA)=\mu(A)\) for every \(g\in G\) and every \(A\subseteq G\). The original statement says that a negative answer implies that \(F\) is not amenable. Its supplied \emph{Comment of 2026} additionally asserts that a positive answer implies that \(F\) is amenable, making this problem equivalent to Notebook Problem 12.20. The comment attributes this converse to an Ore-condition/amenability equivalence for groups satisfying Kaplansky's zero-divisor conjecture, and asserts that \(F\) satisfies that conjecture. Here the relevant zero-divisor assertion is that a product of two nonzero elements of \(k[F]\) cannot be zero. These implications and the zero-divisor assertion are retained as source-reported qualifications, not independently verified results. The text of Problem 12.20 and the cited preprint were not supplied, so their precise formulations and the scope of the cited equivalence remain unverified.
\subsection{Short English statement}
For an arbitrary field \(k\), must every two elements \(a,b\) of the group algebra of Thompson's group \(F\) admit multipliers \(u,v\), at least one nonzero, with \(au=bv\)? The supplied 2026 comment reports that this question is equivalent to amenability of \(F\), using a zero-divisor qualification and a cited result of D. Kielak. The registry labels the problem open; this draft does not independently verify that status or the reported equivalence.
\subsection{Sources}
\begin{itemize}
\item Local UnsolvedMath registry, supplied complete source record: permanent ID \(701130\), exact source code \texttt{catalogue entry}, title \emph{Kourovka Notebook Problem 15.42}. Upstream revision: \texttt{d371ee917ac9647fbb06cd6fc529deec6f3dda82}. The record was updated on \texttt{2026-10-16T14:00:00Z} and supplies the status \texttt{open}; these are registry metadata, not independent mathematical verification.
\item \emph{The Kourovka Notebook (Unsolved Problems in Group Theory)}, Problem 15.42, proposed by V. S. Guba, 15th issue (2002); supplied source also identifies the 21st issue edition (2026), \texttt{arXiv:1401.0300v48}. The collection description names E. I. Khukhro and V. D. Mazurov as editors. The statement refers to Notebook Problem 12.20.
\item The supplied \emph{Comment of 2026} cites D. Kielak, a preprint described in the record as dating from 2026, \url{https://arxiv.org/pdf/1605.09133v2}, for the asserted equivalence under Kaplansky's zero-divisor qualification. This is supplied dated commentary, not an independently verified literature assessment; the record provides no preprint title or separate literature-verification date.
\end{itemize}
\end{document}
